\section{Introduction}
\label{sec:intro}

% Mixed precision
Mixed-precision weights let an accelerator serve large language models (LLMs) with far fewer weight bits and little quality loss: a Fisher-guided P4--P8 posit allocation of Qwen2.5-14B at 5.5 bits per weight keeps WikiText-2 perplexity at 6.11, against 5.88 for FP16 and 6.41 for uniform 5-bit weights.
On a bit-serial datapath such as FlexPosit~\cite{gao2026flexposit}, a weight tile takes 94 cycles at P4 and 174 at P8, so the compute time $C$ of a block follows the precisions of its tiles while its HBM time $M$ follows its weight bytes.

% Energy is not handled well: three problems
Serving, however, is increasingly limited by energy rather than by peak throughput~\cite{iea2025energyai,tpuv4:isca23}, and precision alone handles energy poorly, for three reasons.
First, precision reaches only part of the energy: under power and clock gating at peak voltage and frequency (V/$f$), the MAC arrays take 35--41\% of the energy of a Qwen2.5-14B request from P4 to P8, and the rest goes to memory logic, HBM and vector units that respond differently or not at all.
Second, time saved is not turned into energy saved: a low-precision tile finishes early but the whole core stays at one V/$f$, so arrays wait for HBM while logic stays powered, and memory logic alone takes 25--32\% of the energy.
Third, energy management does not see precision: dynamic voltage and frequency scaling (DVFS) for accelerators can spend the slack under the time-to-first-token (TTFT) and time-per-output-token (TPOT) objectives of each request~\cite{ascend_dvfs:asplos25,xue2026finedvfs}, but it was built for fixed precision, and costing tiles at the worst-case precision loses up to 46\% of the attainable saving at the same objectives.

% Insight
Our insight is that the components want different voltages at the same time, and that this holds per block: under a 5\% TTFT budget the vector rail stays near 0.69\,V at every precision while the MAC rail settles at 0.55--0.59\,V and array logic near 0.47\,V. Precision varies between blocks, not within them: 76\% of its variance lies between blocks, and a voltage change costs a fifth or more of a tile but at most about 1\% of a block, and splitting blocks into sub-blocks gains nothing on our workloads.
A block is therefore the right unit, and since precision is fixed at compile time, the precision map tells, before a request runs, how far each block's $C$ lies from its $M$.
We use this to minimize energy on production traces while requests meet latency targets taken from MLPerf Inference~\cite{mlperf_llama31,mlperf_llama270b}.

% Three challenges; C2 comes from the answer to C1
Acting on this insight raises three challenges (\Cref{sec:motivation}).
\emph{C1, precision-dependent imbalance}: precision changes each component's load by a different amount, so components need their own operating points; a single rail per core cannot provide this, and switched per layer or per phase it saves 0.8--3.9\% over power and clock gating.
\emph{C2, crossing latency}: solving C1 with separate domains creates a new cost, because precision shortens matrix tiles but not the dependent vector chains between them, and the domains put synchronizers on these chains; without local forwarding, TTFT grows by 41--46\% and the objectives can no longer be met, and forwarding through shadow registers that alias frontend registers~\cite{xue2026finedvfs} sends every result back across the crossing, which stalls vector kernels by 45--50\% behind a 4-entry return queue and costs 1.4--4.2 points of saving.
\emph{C3, balance shift}: precision decides whether a block is compute-bound ($C>M$) or HBM-bound, so a planner must follow the side each block falls on, and it must do so fast enough to run per request without searching the schedule.

% Our approach: one technique per challenge
We present \sys{} (MIxed-precision accelerator with Domain-level, precision-Aware Scaling), with one technique per challenge.
\emph{T1} (\Cref{sec:arch}) splits each core into four small V/$f$ rails, for the MAC arrays, array logic, vector and attention units, and memory, each with its own regulator, and a voltage change drains only its target rail at a tile boundary.
\emph{T2} (\Cref{sec:arch}) removes the crossing cost that T1 introduces: since most vector results are consumed only inside the vector unit, the unit keeps them in local shadow slots and returns only results the compiler marks as exported, which needs a 4-entry return queue and 42\% less interface area than writing every result back with no stall.
\emph{T3} (\Cref{sec:pivot}) is a causal, search-free controller that reads the precision map, gives each block its free point, the lowest point at which $C$ stays within $M$, and at dispatch spends the slack it measures on the steps with the largest energy saved per unit of delay first, with separate TTFT and TPOT budgets.

% Contributions
This paper makes three contributions:
\begin{itemize}
\item An analysis of DVFS for mixed-precision LLM accelerators that identifies three energy problems of mixed precision and the challenges of precision-dependent imbalance, crossing latency and balance shift, and quantifies how single rails, write-back-all forwarding, and worst-case and average costing fail (\Cref{sec:motivation}).
\item \sys{}, a multi-domain accelerator core with per-domain switching and an export-only vector interface (T1, T2), and a search-free causal controller (T3).
\item An RTL-calibrated request-level evaluation whose P4--P8 matrix timing predicts held-out layers within 1.94\% and held-out tile shapes within 1.47\%, and whose fast planning model predicts the latency of held-out multi-domain plans within 2.3\% of the request replay. On Azure traces with MLPerf-derived targets, \sys{} saves \TBD{x--y\%} of energy over an accelerator at peak V/$f$ without gating, the baseline of prior fine-grained DVFS work~\cite{xue2026finedvfs}, of which power and clock gating account for \TBD{x} points. We also report what does not help: per-block costing about matches average costing, and allocating precision for energy adds at most 0.4 points (\Cref{sec:results}).
\end{itemize}
